\paragraph{Post-test diagnostic controls.}
The primary endpoint and all original outcomes above are retained.
The following controls and support audit were added after inspecting the test;
they are descriptive checks, not a revised primary analysis. The programmed
rule $G_v=E_vp_v$ uses supplied eligibility and root participation and ignores
query occupancy. It requires no training and is a diagnostic, not a learned
label-efficiency competitor.

\begin{table}[htbp]\centering\small
\begin{tabular}{llrr}
\toprule Teacher & Procedure & BA (\%) & Equal-type acc. (\%) \\
\midrule
BE & Native & 84.64 & 51.70 \\
BE & CF & 90.41 & 86.94 \\
BE & Eligibility/participation & 99.71 & 81.27 \\
GM & Native & 50.00 & 35.44 \\
GM & CF & 100.00 & 100.00 \\
GM & Eligibility/participation & 53.03 & 62.34 \\
\bottomrule
\end{tabular}
\caption{Same final queries and sufficient-type key; means over ten seeds.
The programmed diagnostic ignores occupancy. BA averages sensitivity and
specificity; equal-type accuracy first averages within each distinct
type, then across types within seed. The diagnostic is post-test.}
\label{tab:prospective-controls}
\end{table}

For BE, the programmed rule scores 99.71\% BA, above CF's 90.41\%,
whereas CF scores 86.94\% equal-type accuracy versus 81.27\%.
A single negative type accounts for 95.05--96.30\% of each 4,000-query bank.
Class balancing therefore leaves the negative subtypes highly imbalanced;
high occurrence-level BA alone does not identify the occupancy condition.
Among eligible active BE queries, native and the programmed rule accept
995/995 free options but reject 0/228 occupied options.
CF rejects 228/228 occupied options but accepts only
814/995 free options. These are pooled descriptive counts over
ten banks, not independent training replications. Appendix~\ref{app:prospective-audit}
reports the full confusion strata.

GM's native control receives 4,344--4,354 positive COMMIT labels and zero
negative COMMIT labels in ASSIGN/REPAIR per seed, while CF receives
1,038--1,098 negatives and the test has 2,420--2,491 negatives per seed.
The native model accepts all test queries, giving BA 50\%; the contrast
therefore demonstrates learning rejections absent from native color supervision.
DETECT has separate positive/negative conflict supervision and is excluded
from these color-label counts. This result does not compare CF with another
acquisition strategy supplying informative negative confirmations.

\paragraph{Support of the held-out queries.}
There are 17--23 distinct sufficient test types per BE seed and 82--92
per GM seed, despite 4,000 occurrences. For these color oracles
$O_T=E_vp_v(1-o_v)$, with $o_v$ the queried-option occupancy. A coarser
configuration $(\text{phase},E_v,p_v,o_v)$ is therefore oracle-sufficient.
CF training has already seen configurations covering 99.475--100.000\% of BE
test occurrences and 100\% of GM test occurrences. The original fine-key
holdout remains disjoint: novelty is defined by contextual combinations
rather than by this coarser configuration. There is no separate evaluation
restricted to novel logical conditions. These tests do not establish extrapolation to new logical
conditions. Per-seed support counts and logical overlap are in
Appendix~\ref{app:prospective-audit}; no split, label or model was changed.
